\documentclass[preview]{standalone}

\usepackage{amsmath}
\usepackage{amssymb}
\usepackage{stellar}
\usepackage{definitions}

\begin{document}

\id{hausdorff-measure-introduction}
\genpage

\section{Hausdorff measure}

\begin{snippet}{hausdorff-measure-intuition}
    Lebesgue measure is suitable for measuring sets that have the same
    dimension as the ambient space. For example, in \(\realnumbers^2\)
    it measures areas, while in \(\realnumbers^3\) it measures volumes.

    However, many geometrically interesting sets have lower dimension
    than the space in which they are embedded. A simple plane curve,
    seen as a subset of \(\realnumbers^2\), has zero area:
    \[
        \lambda^2(\gamma)=0.
    \]
    This does not mean, however, that the curve is ``negligible'' from a
    geometric point of view: it may have positive length.

    The idea behind Hausdorff measure is therefore to measure a set using
    the correct dimension. For this reason, we introduce a family of measures
    \[
        \mathcal H^s,
        \qquad s\geq 0,
    \]
    where the parameter \(s\) represents the dimension with which we want to
    observe the set.

    For example:
    \[
        \mathcal H^1
        \quad\text{measures lengths,}
        \quad
        \mathcal H^2
        \quad\text{measures areas,}
        \quad
        \mathcal H^3
        \quad\text{measures volumes.}
    \]
    When \(s=n\), the Hausdorff measure \(\mathcal H^n\) coincides, up to a
    normalization constant, with the \(n\)-dimensional Lebesgue measure on
    \(\realnumbers^n\).

    The main advantage is that the parameter \(s\) does not have to be an
    integer. This makes it possible to describe sets whose geometric dimension
    lies between two integer dimensions, as happens for many fractal sets.
\end{snippet}

\begin{snippetdefinition}{hausdorff-measure-definition}{Hausdorff measure}
    Let \(E \subseteq \realnumbers^n\) and let \(s \geq 0\). For every
    \(\delta > 0\), define
    \[
        \mathcal H^s_\delta(E)
        =
        \inf\left\{
            \sum_{i=1}^\infty {(\operatorname{diam} U_i)}^s
            \suchthat
            E \subseteq \bigcup_{i=1}^\infty U_i,
            \operatorname{diam} U_i \leq \delta
        \right\}.
    \]
    The \emph{\(s\)-dimensional Hausdorff measure of \(E\)} is defined as
    \[
        \mathcal H^s(E)
        \triangleq
        \lim_{\delta \to 0^+} \mathcal H^s_\delta(E)
        =
        \sup_{\delta > 0} \mathcal H^s_\delta(E).
    \]
    Here, the infimum is taken over all countable covers of \(E\) by sets
    of diameter at most \(\delta\).
\end{snippetdefinition}

\end{document}
